\chapter{Evaluation Plan}
\label{ch:experiments}

The evaluation has two levels (Table~\ref{tab:exp-overview}). Policies that do
not change what a model is shown are compared by \emph{replay}: sessions are
recorded once, and every policy is charged from the same recorded candidates.
Handovers that do change what the strong model is shown are run live.

\begin{table}[htbp]
\centering
\caption{The five experiments.}
\label{tab:exp-overview}
\footnotesize
\begin{tabularx}{\textwidth}{@{}>{\raggedright\arraybackslash}p{0.25\textwidth} >{\raggedright\arraybackslash}p{0.06\textwidth} >{\raggedright\arraybackslash}X@{}}
\toprule
Experiment & RQ & What is run\\
\midrule
E0 the warm substrate & RQ2 & every distinct candidate of the study validated twice, once on the cold path and once on M0: cost per step and any disagreement in verdict; no model call\\
E1 single-tier sessions, live & RQ1 & each model alone on the $100$ evaluation bugs, $3$ sessions per bug: the cheap model to $20$ candidates, the strong model to $10$ rounds\\
E2 escalation policies, replayed & RQ3 & six policies charged from the recorded sessions of E1; no model call\\
E3 handover, live & RQ4 & on $50$ bugs, the strong tier restarted with each of three notes in the sessions that step up\\
E4 whole system, live & RQ5 & the full bot against the strong model alone on the other $50$ bugs\\
\bottomrule
\end{tabularx}
\end{table}

\section{Bugs and Setup}
\label{sec:exp-selection}

\paragraph{Bugs.} Defects4J~\citep{just2014defects4j} version 3.0.1 has $854$
real Java bugs from $17$ projects. A bug is used if (i) it can be reproduced,
that is, its failing tests fail before the fix and pass after it; (ii) the fix
changes one source file, and in it only the bodies of existing methods; and
(iii) its test suite has at least ten test classes and runs in under $5$
minutes. From these, $120$ bugs are drawn at random, spread evenly over the
projects. Twenty form a \emph{development set}, used to check the pipeline and
fix every setting in advance. One hundred form the \emph{evaluation set}, on
which nothing is tuned; it is split at random into two halves of $50$ for E3
and E4.

\paragraph{Faulty code given.} Finding the faulty code is not studied. As is
common in LLM-based repair~\citep{xia2024chatrepair}, the model is given the
methods that the developer's fix changes, and every candidate replaces these
methods and nothing else. This favours repair success~\citep{liu2019flbias}
but is the same for every configuration.

\paragraph{Model pair.} Two tiers of one family through an API: GPT-6 Luna as
the cheap tier and GPT-6 Sol as the strong tier, at the dated snapshots in use
when the runs start. Both models reason before they answer, by default at
effort \emph{medium}, and reasoning tokens are billed as
output~\citep{openai2026models}. The strong tier keeps that default. The cheap
tier runs with reasoning off: its part is to supply many cheap, independent
candidates, and reasoning would raise the price of each one and slow the
tier that every session waits for first. Whether this costs the cheap tier too
many repairs is checked on the development set
(Table~\ref{tab:exp-provisional}). Defects4J is likely memorised by current
models~\citep{ramos2024memorizing}. If the cheap tier repaired nearly every
bug, or the strong tier repaired every remaining bug at once, there would be
little to measure. The pair is therefore confirmed on the development set by a
rule fixed in advance. The cheap tier must repair between $20\,\%$ and
$80\,\%$ of the development bugs within its budget, and the strong tier must
need more than one round on at least a third of the bugs the cheap tier fails.
If the cheap tier repairs more than $80\,\%$, or the strong tier fails the
second condition, the pair moves down one step: a local 7B model
(\texttt{qwen2.5-coder:7b}) becomes the cheap tier and Luna the strong one. If
the cheap tier repairs less than $20\,\%$, the pair is kept and the result is
read as in Chapter~\ref{ch:outcomes}. The thesis reports which pair was used
and why.

\paragraph{Unit of testing.} Tests are scheduled by test class. A candidate is
plausible when it passes every scheduled class.

\section{What Is Compared}
\label{sec:exp-arms}

Six policies are compared in E2 (Table~\ref{tab:exp-arms}). All of them use
the same validation layer, so they differ only in which model is asked.

\begin{table}[htbp]
\centering
\caption{The policies compared.}
\label{tab:exp-arms}
\footnotesize
\begin{tabularx}{\textwidth}{@{}>{\raggedright\arraybackslash}p{0.15\textwidth} >{\raggedright\arraybackslash}X >{\raggedright\arraybackslash}p{0.12\textwidth}@{}}
\toprule
Policy & Rule & Role\\
\midrule
Cheap-only & the cheap model, up to $20$ candidates & reference\\
Strong-only & the strong model, up to $10$ rounds, conversational & baseline\\
Fixed-5 & step up after $5$ cheap candidates & baseline\\
Fixed-10 & step up after $10$ cheap candidates & baseline\\
\Step & Algorithm~\ref{alg:step}: step up after a stalled call, with at most $20$ cheap candidates & proposed\\
Hindsight & for each session, the switch point that costs least, chosen with knowledge of the outcome & bound\\
\bottomrule
\end{tabularx}
\end{table}

\paragraph{Hindsight} is not a policy that can be deployed. For each session
it picks, among the switch points of $0$, $5$, $10$, $15$ and $20$ cheap
candidates, the one whose session costs least in money, and it prefers a
switch point that ends with a plausible patch to one that does not. It
therefore marks the most that any router could save by choosing when to
switch, the learned routers of Section~\ref{sec:rel-f} included. The distance
between \Step\ and Hindsight is the room left for a better trigger.

\paragraph{Two readings of the budget.} In the main reading, every escalating
policy gives the strong tier its full $10$ rounds, so
Proposition~\ref{prop:nolost} applies and no repair is lost. In a second
reading the total number of attempts is held at $10$: five cheap candidates,
then five strong rounds. There the repair rate can fall, and it is tested.

\paragraph{Handovers (RQ4).} Cold, Evidence, Evidence and patches, and
Transcript (Section~\ref{sec:sys-controller}), each with \Step.

\paragraph{Variants of \Step.} A stricter stall, in which a call has stalled
only when none of its candidates makes progress, is compared by replay. Three
variants change what the cheap model is asked, so they are run live, on the
development set only: other numbers of candidates per call ($1$, $3$ and
$10$), a conversational cheap tier instead of blind re-sampling, and a cheap
tier that reasons at effort \emph{low}.

\section{How Cost Is Measured}
\label{sec:exp-measure}

\paragraph{Audit.} Every distinct candidate is compiled once and run on every
scheduled test class, on the cold path and on the warm path, recording the
result and time of each class. From this record the validation cost of any
policy on that candidate is computed: compilation, start-up where it applies,
and the classes run up to the first failure, in the order the session would
have used. All policies are thus charged from the same measurement. The
verdict of record is the one of the cold path.

\paragraph{How time is measured.} On the cold path the classes of a candidate
run one after another in a JVM started for that candidate, as build tools run
a test suite; on the warm path they run in the live runtime. The runner
records for each class its result, its wall-clock time and the CPU time of the
JVM process while the class runs. A class that runs longer than $60$ seconds
counts as failed. The main unit for validation is CPU seconds, which depend
less on the load of the machine than wall-clock time; wall-clock time is
recorded as well. A session on the warm path is also charged, once, for
starting the runtime and for its calibration run.

\paragraph{What is recorded.} For each model call: tier, latency and tokens.
For each audit and each path: compile time, start-up, and the result,
wall-clock time and CPU time of each test class. For each model call the
record also holds the reasoning tokens and the input tokens that the API
reports as read from or written to its prompt cache. Costs are reported in
time, CPU time, tokens and money, using one fixed price
list~\citep{openai2026models}: \$0.10 and \$0.50 per million input and output
tokens for Luna, \$2 and \$10 for Sol, and \$0.05 per CPU-hour of validation.
A local model has no token price. It is charged in time, and in a sensitivity
analysis at a stated price per GPU-hour.

\paragraph{Prompt caching.} The API caches the prefix of any prompt of at
least $1{,}024$ tokens by default. A token read from the cache costs a tenth
of the input price, and a token written to it $1.25$
times~\citep{openai2026caching}. What a request finds in the cache depends on
what was sent in the minutes before it, so the bill of a recorded call is not
a property of its session, and a replayed policy cannot be priced from it.
Money is therefore reported in two readings that depend on the session
alone. In the first, every
input token pays the input price. In the second, a cache holds, within one
session, every prefix that a tier has already sent: a blind call re-reads its
prompt, and a round of a conversation reads the previous prompt and writes
what was added. The first reading ignores the cache; the second gives a
session the most it can get from the cache on its own. The cache is expected
to favour the strong tier, whose conversation sends its history again in
every round. The conclusions of RQ3 and RQ5 are drawn from whichever reading
shows the smaller saving, and the tokens that the API reported as cached are
shown next to the model.

\paragraph{Reliability checks.} Each test class runs three times on the fixed
program; classes that change outcome are excluded as flaky. On a random tenth
of the candidates, the audit is repeated in two other orders, to check that
results do not depend on test order. Every candidate that passes the scheduled
classes is run once on the whole suite: a failure there would mean that the
skipped classes hid it, so the candidate is counted as rejected and the case
is reported. That run is the CI run any proposed patch receives; it is the
same under every policy and is not charged to the session.

\section{E0 and E1: Single-Tier Sessions, Audited}
\label{sec:exp-e1}

\paragraph{The loops.} The cheap tier draws candidates by blind re-sampling:
each call sends the first prompt and returns $n = 5$ candidates, at
temperature $1.0$ and with reasoning off. The strong tier holds a conversation
in the style of ChatRepair~\citep{xia2024chatrepair}, one candidate per round,
with the feedback of Table~\ref{tab:exp-prompts}, at its default reasoning
effort and sampling. The
first prompt is the same for both tiers and gives the faulty methods, a
failing test and its error message (Figure~\ref{fig:exp-prompt}). A session
ends at the first plausible patch or when its budget is spent. If an endpoint
cannot return several candidates for one request, the five are drawn by five
requests; the prompt is then paid five times, and the break-even is reported
at that price.

\begin{table}[htbp]
\centering
\caption{Feedback to the strong tier after each round.}
\label{tab:exp-prompts}
\footnotesize
\begin{tabularx}{\textwidth}{@{}>{\raggedright\arraybackslash}p{0.27\textwidth} >{\raggedright\arraybackslash}X@{}}
\toprule
Outcome & Feedback to the model\\
\midrule
A test fails & the first failing test, its line and error; ask for a different fix\\
Does not compile & the compiler error; ask for a corrected version\\
Unreadable answer & the answer could not be read as the given methods\\
Repeat of its own candidate & the patch was already tried, and which test rejected it\\
\bottomrule
\end{tabularx}
\end{table}

\begin{figure}[htbp]
\centering
\begin{minipage}{0.95\textwidth}
\begin{Verbatim}[frame=single,rulecolor=\color{boxline},fontsize=\footnotesize,framesep=2mm]
The following Java method contains a bug:

    private static void modify(Calendar val, int field, boolean round) {
        ...
    }

It makes this test fail:
    Test:  DateUtilsTest::testTruncateLang59
    Error: AssertionFailedError: Truncate Calendar.SECOND
           expected:<Sun Oct 31 01:02:03 MDT 2004> but was:<... MST 2004>

Return the complete fixed method in one Java code block.
\end{Verbatim}
\end{minipage}
\caption{The first prompt for Defects4J bug Lang-65.}
\label{fig:exp-prompt}
\end{figure}

\paragraph{Same draws for every policy.} Model answers are cached under the
full prompt and a draw index (bug, session, tier, call). The index omits the
policy, so two policies that send the same prompt at the same point receive
the same answer. Audits are cached by normalised patch, so each patch is
tested once.

\paragraph{What E1 answers.} The recorded sessions give RQ1: where the cost of
each model goes, split into model calls, building, start-up and tests, and by
how the session ended. Each session is priced twice: on the warm substrate,
and as a plain loop that builds and starts a runtime for every candidate and
validates repeats again. The sessions are also the streams that E2 replays.

\paragraph{What E0 answers.} The audits of the same candidates give RQ2: what
each step of a validation costs on each path, how often the two paths disagree
on a verdict, how many bugs fail the calibration and fall back to the cold
path, and how often a candidate forces a new runtime. They also give what the
schedule leaves to any other order: the test time it spends before the first
failure, against the cheapest rejecting test in hindsight.

\section{E2: Escalation Policies by Replay}
\label{sec:exp-e2}

\paragraph{Why replay is exact.} A policy with a cold handover never changes a
prompt. What a model is sent depends only on its own earlier candidates and on
what is reported about them. The reported failure of a candidate does not
depend on the session (Section~\ref{sec:sys-warm}), and a repeat of the other
tier's candidate is reported as an ordinary failure
(Section~\ref{sec:sys-memory}). The session of such a policy is therefore a
prefix of the recorded cheap session, followed by a prefix of the recorded
strong session of the same bug and session number. Its cost is read from the
records: the calls it would have made, and from the audits the validations it
would have run, in the order its own memory would have scheduled them. Every
policy thus sees exactly the same candidates, and the costs differ only by the
policy. In the implementation a replay is the bot itself, run on recorded
answers only: a policy that needed an answer nobody recorded stops with an
error instead of calling a model.

\paragraph{What is measured.} For each policy: the cost per repaired bug in
money, tokens, time and CPU time; the share of sessions repaired at the cheap
tier; the break-even margin of Proposition~\ref{prop:breakeven}; the time
spent at the cheap tier in sessions that step up; and the distance to
Hindsight.

\paragraph{A check, not a test.} In the main reading a policy repairs every
session that Strong-only repairs. A session in which it does not is a defect
of the harness, and is investigated as one.

\section{E3 and E4: Handover and Whole System, Live}
\label{sec:exp-live}

\paragraph{E3: handover.} On $50$ evaluation bugs, three sessions each, E2
says which sessions step up under \Step. In each of them the strong tier is
started again with each of three handovers: Evidence, Evidence and patches,
and Transcript. The Cold handover is the recorded strong session itself. The
four are compared on the same sessions: rounds of the strong tier until a
plausible patch, repairs within its budget, and tokens. Two more rates are
read from the Evidence Memory. The \emph{refuted-again rate} is the share of
strong candidates that fail with a failure signature the cheap tier had
already produced. The \emph{anchoring rate} is the share of strong candidates
that equal a cheap candidate after normalising. A note that helps should lower
the first without raising the second.

\paragraph{E4: whole system.} The full bot runs on the other $50$ evaluation
bugs, three sessions each: \Step, the handover chosen on E3 by the rule of
Table~\ref{tab:exp-provisional}, and the Budget Guard, which only logs when it
would stop. Its stopping is then replayed, which separates the effect of
stopping from the rest. The bot is compared with Strong-only on the same bugs,
priced in two ways: with the same validation layer, which isolates escalation,
handover and stopping; and as a plain conversational loop, which builds and
starts a runtime for every candidate and validates repeats again. The second
is the comparison of RQ5. The main measure is the cost per repaired bug; the
repair rate must not fall by more than a margin fixed in advance. Choosing the
handover on one half and testing the system on the other keeps the choice from
being tested on the bugs that made it.

\section{Metrics}
\label{sec:exp-metrics}

Table~\ref{tab:exp-metrics} lists the metrics. Costs are summed over sessions
before dividing, and every comparison is made per bug.

\begin{table}[htbp]
\centering
\caption{Metrics.}
\label{tab:exp-metrics}
\footnotesize
\begin{tabularx}{\textwidth}{@{}>{\raggedright\arraybackslash}p{0.25\textwidth} >{\raggedright\arraybackslash}X >{\raggedright\arraybackslash}p{0.07\textwidth}@{}}
\toprule
Metric & Definition & RQ\\
\midrule
Cost per repaired bug & total cost of all sessions divided by the number of bugs repaired; in money (both readings of Section~\ref{sec:exp-measure}), tokens, time and CPU time & 1, 3, 5\\
Cost breakdown & share of each tier's model calls, of building, of start-up and of tests in the total, by how the session ended & 1\\
Validation cost per candidate & on the cold path and on M0; rate of disagreement in verdict; share of bugs that fall back & 2\\
Cheap-tier repair share & share of sessions repaired before any strong call & 3\\
Break-even margin & the cheap tier's repair share under the policy, minus its mean number of candidates times the price ratio $c_1/c_2$ & 3\\
Escalation delay & time spent at the cheap tier in sessions that step up & 3\\
Distance to Hindsight & cost above Hindsight, as a share of the gap between Strong-only and Hindsight & 3\\
Repair rate difference & share of sessions with a plausible patch, against Strong-only & 3, 4, 5\\
Strong rounds to repair & rounds of the strong tier until a plausible patch & 4\\
Refuted-again rate & share of strong candidates that fail with a signature the cheap tier had already produced & 4\\
Anchoring rate & share of strong candidates equal, after normalising, to a cheap candidate & 4\\
Cost of failed sessions & share of the total spent on sessions without a plausible patch & 1, 5\\
Match with the developer's fix & share of plausible patches equal, after normalising, to the developer's fix, per tier & 5\\
\bottomrule
\end{tabularx}
\end{table}

\section{Statistics}
\label{sec:exp-size}

The bug is the unit of analysis; the sessions of a bug are averaged first. The
two main comparisons are \Step\ against Strong-only in cost, and each note
against Cold in strong rounds. They use a paired Wilcoxon test on the per-bug
log-ratio, with a Holm correction, and are reported with bootstrap intervals
over bugs. With $100$ bugs, a saving of $5\,\%$ is detected with power $0.8$
if the standard deviation of the log-ratio is at most $0.15$. If the spread is
larger, the smallest detectable saving is reported. The comparison of the
notes has fewer units: only the bugs of E3 on which a session steps up, at
most $50$. Their number and the smallest difference they can detect are
reported with the result, and a null there is read as an absence of evidence,
not as equivalence.

In the main reading of E2 the repair rate cannot fall, so it is checked rather
than tested (Section~\ref{sec:exp-e2}). In the equal-attempts reading, and in
E3 and E4, the repair rate is tested for non-inferiority with a margin of $10$
percentage points. With $50$ bugs, that margin is reachable if the standard
deviation of the paired difference is at most $0.24$. If the development set
shows more spread than assumed, the number of sessions per bug is raised
before the number of bugs. The assumed values are provisional and checked on
the development set (Table~\ref{tab:exp-provisional}).

\section{Compute and API Budget}
\label{sec:exp-compute}

Most of the compute goes to audits, and most of the money to the strong tier.
Table~\ref{tab:exp-compute} gives the counts. They are upper bounds: a session
stops at the first plausible patch, and policies share cached answers.

\begin{table}[htbp]
\centering
\caption{Calls and audits, as upper bounds.}
\label{tab:exp-compute}
\footnotesize
\begin{tabularx}{\textwidth}{@{}>{\raggedright\arraybackslash}p{0.2\textwidth} >{\raggedright\arraybackslash}X@{}}
\toprule
Work & Upper bound\\
\midrule
Flakiness check & $3$ suite runs per bug\\
E0 & each distinct candidate audited on both paths\\
E1, cheap tier & $100$ bugs $\times$ $3$ sessions $\times$ $4$ calls of $5$ candidates $= 1{,}200$ calls\\
E1, strong tier & $100$ bugs $\times$ $3$ sessions $\times$ $10$ rounds $= 3{,}000$ calls\\
E2 & no model call\\
E3 & $50$ bugs $\times$ $3$ sessions $\times$ $3$ notes $\times$ $10$ rounds $= 4{,}500$ strong calls\\
E4 & $50$ bugs $\times$ $3$ sessions $\times$ $10$ rounds $= 1{,}500$ strong calls; cheap calls shared with E1\\
Audits & one per distinct candidate, shared by all configurations; at most $4$ CPU-hours per bug\\
\bottomrule
\end{tabularx}
\end{table}

The strong tier makes at most $9{,}000$ calls. If every one carried $6{,}000$
input and $3{,}000$ output tokens, reasoning included, they would cost about
\$380 at list prices, and less with the prompt cache;
the $1{,}200$ cheap calls cost under \$2. The development set adds at most
$900$ strong calls ($20$ bugs, three sessions, the $15$ rounds of
Table~\ref{tab:exp-provisional}), about \$40, and its cheap-tier variants a
few dollars. If the estimate from the development set exceeds the budget, cuts
are made in a fixed order: the Transcript handover of
E3 is run on half of its bugs, then the sessions per bug of E3 are lowered
from three to two. If the pair moves down one step, the bound falls by a
factor of twenty.

\section{Pre-Registration}
\label{sec:exp-prereg}

This proposal fixes the bugs, settings, metrics and analyses in advance.
Table~\ref{tab:exp-provisional} gives the values that still depend on data,
each with a provisional value and the rule that fixes it. They are frozen
before E1 starts on the evaluation set, and any later change is reported. The
experiments run in order: E0 and E1, then E2, then E3, then E4.

\begin{table}[htbp]
\centering
\caption{Provisional values and the rules that fix them.}
\label{tab:exp-provisional}
\footnotesize
\begin{tabularx}{\textwidth}{@{}>{\raggedright\arraybackslash}p{0.2\textwidth} >{\raggedright\arraybackslash}p{0.24\textwidth} >{\raggedright\arraybackslash}X@{}}
\toprule
Value & Provisional & Rule\\
\midrule
Model pair & Luna, then Sol & development set: kept if the cheap tier repairs $20\,\%$ to $80\,\%$ of the bugs within its budget and the strong tier needs more than one round on at least a third of the rest; otherwise as in Section~\ref{sec:exp-selection}\\
Candidates per cheap call, $n$ & $5$, at temperature $1.0$ & development set: the largest $n$ of $1$, $3$, $5$ and $10$ at which at least half of the candidates of a call are distinct after normalising\\
Reasoning effort & cheap tier: none; strong tier: medium & development set: none is kept if the cheap tier still repairs at least $20\,\%$ of the bugs within its budget; otherwise low. The strong tier is not changed\\
Cheap budget $B_1$ & $20$ candidates, for \Step\ and for Cheap-only & not changed\\
Stall rule of the cheap tier & a call in which at most one candidate makes progress & development set: kept if it loses at most $5\,\%$ of the cheap tier's repairs; otherwise a call in which no candidate makes progress\\
Strong budget $R_2$ & $10$ rounds & development set, run with $15$ rounds: kept if fewer than $5\,\%$ of the repaired strong sessions need more than $10$; otherwise $15$\\
Stall count $r$ of the strong tier & $3$ rounds without progress & development set: the smallest $r$ whose replay loses at most $5\,\%$ of the plausible sessions\\
Session budget $B$ & $900$ s of wall-clock time & development set: the $90$th percentile of the time to a plausible patch in the Strong-only sessions\\
Handover used in E4 & chosen on E3 & on the $50$ bugs of E3: the fewest strong rounds per repaired bug, with a repair rate within the margin of Cold; otherwise Cold\\
SD of the per-bug log-ratio & $0.15$ & measured; if larger, the smallest detectable saving is reported\\
SD of the repair-rate difference & $0.24$ & measured; if larger, the sessions per bug are raised, not the margin\\
Non-inferiority margin & $10$ percentage points & not changed\\
Bugs & $120$ accepted: $20$ development, $100$ evaluation ($50$ for E3, $50$ for E4) & the share of Defects4J bugs that pass the rules is reported from the selection log\\
\bottomrule
\end{tabularx}
\end{table}
